%! AP Statistics — Unit 9, Lesson 9.4 — Homework (Answer Key)
\documentclass[10pt]{article}
\usepackage{apstats-article}
\usepackage{apstats-key}
\usepackage{grading}
\gradingassignment{apstats-u09-l04-hw}
\gradingcourse{advanced_statistics}

\begin{document}

\pageheader{Unit 9, Lesson 9.4}{Homework --- Answer Key}

\vspace{0.1in}

\textbf{Problem 1.} A random sample of 28 college students was surveyed. The
explanatory variable is hours per week spent on social media ($x$) and the response
variable is GPA ($y$). A residual plot shows random scatter with no curve and constant
spread. The residuals appear approximately normal.

\begin{enumerate}[label=\alph*.]

  \boxguard[15]
  \item Name the appropriate inference procedure.

        \gradedtop{\gradedblank[id=hw1a, type=rubric, points=1, display={$t$-test for slope}]{\linewidth+6pt}{\gradedlinesheight{1}}}

  \boxguard[15]
  \item Define $\beta$ in context and state H$_0$ and H$_a$.

        $\beta = $ \ans{the true slope of the population regression line relating GPA
        to hours per week on social media for college students}\gradedkeyline{\gradedblank[id=hw1b1, type=rubric, points=1, display={\bfseries\boldmath True slope, GPA vs.\ hours/week on social media, all college students.}]{\linewidth+6pt}{\gradedlinesheight{9}}}

        \boxguard[15]
        H$_0$: \ans{$\beta = 0$}\gradedkeyline{\gradedblank[id=hw1b2, type=rubric, points=1, display={\bfseries\boldmath $\beta = 0$}]{\linewidth+6pt}{\gradedlinesheight{9}}}

        \boxguard[15]
        H$_a$: \ans{$\beta \ne 0$}\gradedkeyline{\gradedblank[id=hw1b3, type=rubric, points=1, display={\bfseries\boldmath $\beta \ne 0$}]{\linewidth+6pt}{\gradedlinesheight{9}}}

  \boxguard[15]
  \item Verify each LINER condition:

        \gradedtop{\gradedblank[id=hw1c, type=rubric, points=3, display={L: Met --- residual plot shows random scatter with no curve. I: Met --- random sample stated; need to assume the college population is at least $10 \times 28 = 280$ students (10\% condition). N: Met --- residuals approximately normal. E: Met --- residual plot shows constant spread. R: Met --- random sample stated.}]{\linewidth+6pt}{\gradedlinesheight{5}}}

  \boxguard[15]
  \item \gradedtop{\gradedblank[id=hw1d, type=rubric, points=2, display={The 10\% condition for I cannot be fully verified without knowing the total size of the college student population. If the college has fewer than 280 students, the condition would not be met.}]{\linewidth+6pt}{\gradedlinesheight{2}}}

\end{enumerate}

\newpage

\textbf{Problem 2.}

\begin{enumerate}[label=\alph*.]

  \boxguard[15]
  \item H$_0$: \ans{$\beta = 0$ (rainfall has no linear effect on umbrella sales)}\gradedkeyline{\gradedblank[id=hw2a1, type=rubric, points=1, display={\bfseries\boldmath $\beta = 0$}]{\linewidth+6pt}{\gradedlinesheight{9}}}

        \boxguard[15]
        H$_a$: \ans{$\beta \ne 0$ (rainfall is linearly related to umbrella sales)}\gradedkeyline{\gradedblank[id=hw2a2, type=rubric, points=1, display={\bfseries\boldmath $\beta \ne 0$}]{\linewidth+6pt}{\gradedlinesheight{9}}}

        \boxguard[15]
        $\beta = 0$ means:\gradedkeyline{\gradedblank[id=hw2a3, type=rubric, points=1, display={daily rainfall has no linear effect on umbrella sales in the population.}]{\linewidth+6pt}{\gradedlinesheight{9}}}

  \newpage
  \boxguard[15]
  \item H$_0$: \ans{$\beta = 0$}\gradedkeyline{\gradedblank[id=hw2b1, type=rubric, points=1, display={\bfseries\boldmath $\beta = 0$}]{\linewidth+6pt}{\gradedlinesheight{9}}}

        \boxguard[15]
        H$_a$: \ans{$\beta > 0$ (one-sided: older runners take longer)}
        (also accept $\beta \ne 0$ as a two-sided alternative)\gradedkeyline{\gradedblank[id=hw2b2, type=rubric, points=1, display={\bfseries\boldmath $\beta > 0$}]{\linewidth+6pt}{\gradedlinesheight{9}}}

        \boxguard[15]
        $\beta = 0$ means:\gradedkeyline{\gradedblank[id=hw2b3, type=rubric, points=1, display={age has no linear effect on marathon finish time in the population.}]{\linewidth+6pt}{\gradedlinesheight{9}}}

  \newpage
  \boxguard[15]
  \item H$_0$: \ans{$\beta = 0$}\gradedkeyline{\gradedblank[id=hw2c1, type=rubric, points=1, display={\bfseries\boldmath $\beta = 0$}]{\linewidth+6pt}{\gradedlinesheight{9}}}

        \boxguard[15]
        H$_a$: \ans{$\beta < 0$ (one-sided: larger classes score lower)}\gradedkeyline{\gradedblank[id=hw2c2, type=rubric, points=1, display={\bfseries\boldmath $\beta < 0$}]{\linewidth+6pt}{\gradedlinesheight{9}}}

        \boxguard[15]
        $\beta = 0$ means:\gradedkeyline{\gradedblank[id=hw2c3, type=rubric, points=1, display={class size has no linear effect on average exam score in the population.}]{\linewidth+6pt}{\gradedlinesheight{9}}}

\end{enumerate}

\begin{teachernote}
For Problem 2b, the question says the researcher ``considers'' a one-sided test.
Accept either $\beta > 0$ (one-sided) or $\beta \ne 0$ (two-sided). The one-sided
alternative is appropriate only if the researcher committed to it before seeing the
data. For Problem 2c, note that the teacher's claim involves $\beta < -1$, but AP
convention tests H$_0$: $\beta = 0$; accept H$_a$: $\beta < 0$ as a valid simplification.
\end{teachernote}

\end{document}
